\documentclass{article}
\usepackage{graphicx}

\date{\today}
\author{Tarik Atlaoui \\ Nicolas Peugnet \\ Kimmeng Ly \\ Max Eliet}

\begin{document}


\begin{titlepage}
	\enlargethispage{2cm}
	\newcommand{\HRule}{\rule{\linewidth}{0.5mm}}
	\center
	\textsc{\LARGE
	Pre-rapport du PSAR 
	} \\[1cm]
	\HRule \\[0.4cm]
	{ \huge \bfseries API générique pour le développement d'applications réparties \\[0.15cm] }
	\HRule \\[4cm]
	\large{Tarik Atlaoui \\[3mm] Nicolas Peugnet \\[3mm] Kimmeng Ly \\[3mm] Max Eliet} \\[3cm]
	09 Mars 2020 \\[3cm]
	\hfill \includegraphics[width=5cm]{logoSU.jpg}
\end{titlepage}

	\newpage
	\pagenumbering{arabic}
		\section{Introduction}
			\subsection{Les applications réparties : qu'est-ce ?}
			 \subsection{Les difficultés à programmer une application répartie}

		\section{Méthode de développement}
			\subsection{API réelle : avantages/inconvénients + exemple API(MPI ou autre)}
			\subsection{API simulation à évenements discrets : qu'est-ce ? + avantages/inconvénients + exemple(PeerSim ou autre)}
				\large{Qu'entend-on par \textit{simulation à événements discrets}? C'est une simulation dont le temps évolue seulement lorsqu'un événement survient sur un noeud, et donc de même l'état du système ne peut être modifié qu'à ces moments là.
\newline \indent On distingue donc deux entités : les noeuds, et les événements qui sont caractérisés par une date de délivrance, un noeud destinataire, et des données.
\newline \indent Les principaux avantages d'un tel type de simulation sont : son déterminisme et donc une capacité à reproduire des bugs, et une charge de calcul réduite aux événements qu'on décide de simuler. Toutefois, il faut savoir trouver un équilibre entre une simulation trop simpliste et une simulation trop précise ralentie par trop d'événements.
\newline \indent Dans notre cas, nous nous sommes dirigés vers PeerSim comme simulateur à événements discrets, car il est codé en Java et possède une API relativement simple d'utilisation. }
			\subsection{Aperçu de l'implémentation d'un anneau avec les deux API}
		
		\section{Motivation et objectif global}
			\subsection{API générique : avantages, modèle de programmation(ici événementiel)}
			\subsection{Code précédent en version générique}
		
		\section{Les étapes de réalisation}
			\subsection{Définition API générique}
			\subsection{Définition des primitives offertes et explication}
			\subsection{Implantation vers MPI}
			\subsection{Implantation vers PeerSim}
			\subsection{API pour un scenario}
			\subsection{Rédaction du rapport}

		\section{Plan de validation}
			\subsection{Comment montrer que cela fonctionne bien, tests unitaires, respect du cahier des charges}
		\section{Planning des tâches}
			\subsection{Une deadline pour chaque tâche dans un ordre chronologique}
\end{document}
